\appendix
\section{Pixel-level MSTNet pipeline}
\label{app:pixel-pipeline}

Figure~\ref{fig:data-flow} in the main text shows the municipal-mean path used for training and for the municipal-mean LOO metrics in Section~\ref{sec:experiments}.
Figure~\ref{fig:data-flow-pixel} diagrams the pixel-wise alternative used after training: soy-masked Sentinel-2 pixels are encoded one by one, climate is broadcast across pixels, soil enters by late fusion, and the mean of the pixel predictions $\hat{y}_p$ is compared with PAM~\eqref{eq:pixel-mean-agg}.
We do not train with that path; it is much slower than municipal-mean training while targeting the same labels.
The diagram is still useful because the post-training evaluation in Section~\ref{sec:results} follows the same forward path without backpropagating through every pixel, and the resulting field $\{\hat{y}_p\}$ can be kept as a map or averaged over any polygon.

\begin{figure}[!htbp]
\centering
\resizebox{\textwidth}{!}{%
\begin{tikzpicture}[
  font=\footnotesize,
  line join=round,
  x=1cm, y=1cm,
  >={Stealth[length=1.8mm, width=1.4mm]},
  shorten >=2.2mm,
  shorten <=0.4mm,
  src/.style={
    draw=black!40, fill=white, rounded corners=2pt,
    align=center, inner xsep=5pt, inner ysep=5pt,
    minimum width=2.35cm, minimum height=1.15cm, line width=0.45pt
  },
  mid/.style={
    draw=black!40, fill=black!2, rounded corners=2pt,
    align=center, inner xsep=5pt, inner ysep=5pt,
    minimum width=2.45cm, minimum height=1.05cm, line width=0.45pt
  },
  enc/.style={
    draw=#1!55!black, fill=white, rounded corners=2pt,
    align=center, inner xsep=6pt, inner ysep=5pt,
    minimum width=2.5cm, minimum height=0.95cm, line width=0.55pt
  },
  dec/.style={
    draw=black!50, fill=black!4, rounded corners=2pt,
    align=center, inner xsep=8pt, inner ysep=6pt,
    minimum width=2.9cm, minimum height=0.95cm, line width=0.55pt
  },
  rowlab/.style={font=\scriptsize\color{black!50}, anchor=east},
  arr/.style={->, line width=0.55pt, draw=black!45},
]

  \node[rowlab] at (-8.0, 0) {Sources};

  \node[src] (s2) at (-6.3, 0) {\scriptsize Sentinel-2};
  \begin{scope}[on background layer]
    \node[src, fill=black!7, draw=black!26, line width=0.4pt,
          inner sep=0pt, fit=(s2), xshift=4pt, yshift=-4pt] {};
    \node[src, fill=black!3, draw=black!30, line width=0.4pt,
          inner sep=0pt, fit=(s2), xshift=2pt, yshift=-2pt] {};
  \end{scope}

  \node[src] (soy) at (-3.15, 0) {%
    \begin{tabular}{@{}c@{}}
      \scriptsize MapBiomas\\[0pt]
      \scriptsize soy mask
    \end{tabular}};

  \node[src] (xav) at (0, 0) {%
    \begin{tabular}{@{}c@{}}
      \scriptsize Xavier\\[0pt]
      \scriptsize weather
    \end{tabular}};
  \begin{scope}[on background layer]
    \node[src, fill=black!7, draw=black!26, line width=0.4pt,
          inner sep=0pt, fit=(xav), xshift=4pt, yshift=-4pt] {};
    \node[src, fill=black!3, draw=black!30, line width=0.4pt,
          inner sep=0pt, fit=(xav), xshift=2pt, yshift=-2pt] {};
  \end{scope}

  \node[src] (soil) at (3.15, 0) {%
    \begin{tabular}{@{}c@{}}
      \scriptsize MapBiomas\\[0pt]
      \scriptsize Solo
    \end{tabular}};

  \node[src] (pam) at (6.3, 0) {\scriptsize IBGE PAM};

  \node[rowlab] at (-8.0, -3.0) {Pixels};

  \node[mid] (p_s2) at (-4.725, -3.0) {\scriptsize Soy pixel series};
  \begin{scope}[on background layer]
    \node[mid, fill=black!7, draw=black!26, line width=0.4pt,
          inner sep=0pt, fit=(p_s2), xshift=4pt, yshift=-4pt] {};
    \node[mid, fill=black!3, draw=black!30, line width=0.4pt,
          inner sep=0pt, fit=(p_s2), xshift=2pt, yshift=-2pt] {};
  \end{scope}

  \node[mid] (p_cl) at (0, -3.0) {\scriptsize Climate\\[-1pt]\scriptsize (broadcast)};
  \node[mid] (p_so) at (3.15, -3.0) {\scriptsize Soil\\[-1pt]\scriptsize (per pixel)};
  \node[mid] (p_lb) at (6.3, -3.0) {\scriptsize Labels};

  \draw[arr] (s2.south) -- ++(0,-0.55) -| (p_s2.north);
  \draw[arr] (soy.south) -- ++(0,-0.55) -| (p_s2.north);
  \draw[arr] (xav.south) -- (p_cl.north);
  \draw[arr] (soil.south) -- (p_so.north);
  \draw[arr] (pam.south) -- (p_lb.north);

  \node[rowlab] at (-8.0, -6.2) {Encoders};
  \node[enc=blue!55!black] (e_s2) at (-4.725, -6.2) {\scriptsize Spectral encoder};
  \node[enc=orange!65!black] (e_cl) at (0, -6.2) {\scriptsize Climate encoder};

  \begin{scope}[on background layer]
    \node[
      fit=(e_s2)(e_cl),
      fill=blue!4,
      draw=blue!25!black!40,
      rounded corners=3pt,
      inner xsep=18pt,
      inner ysep=16pt,
      line width=0.5pt
    ] (encbox) {};
  \end{scope}
  \node[font=\scriptsize\color{blue!40!black!60}, anchor=south]
    at ([yshift=4pt]encbox.north) {MSTNet};

  \draw[arr] (p_s2.south) -- (e_s2.north);
  \draw[arr] (p_cl.south) -- (e_cl.north);

  \node[rowlab] at (-8.0, -9.0) {Decoder};
  \node[dec, minimum width=3.6cm, text width=3.4cm] (dec) at (-2.3625, -9.0)
    {\scriptsize Late fusion + MLP $\Rightarrow\hat{y}_p$};

  \draw[arr] (e_s2.south) -- ++(0,-0.7) -| (dec.150);
  \draw[arr] (e_cl.south) -- ++(0,-0.7) -| (dec.30);
  \draw[arr] (p_so.south) -- (p_so.south |- dec.east) -- (dec.east);

  \node[rowlab] at (-8.0, -11.2) {Aggregate};
  \node[dec] (agg) at (-2.3625, -11.2) {\scriptsize Mean of pixel predictions};

  \draw[arr] (dec.south) -- (agg.north);

  \node[rowlab] at (-8.0, -13.4) {Output};
  \node[dec] (out) at (-2.3625, -13.4) {$\hat{y}_{m,y}$};

  \draw[arr] (agg.south) -- (out.north);
  \draw[arr] (p_lb.south) -- (p_lb.south |- out.east) -- (out.east);
\end{tikzpicture}%
}
\caption{Pixel-wise MSTNet data pipeline used for post-training evaluation and sub-municipal maps, not for training.
Soy-masked Sentinel-2 pixels are encoded individually; municipal Xavier climate is broadcast to pixels and soil enters by late fusion.
The mean of pixel predictions $\hat{y}_p$ yields $\hat{y}_{m,y}$ for comparison with IBGE PAM.}
\label{fig:data-flow-pixel}
\end{figure}

The network weights are shared with the municipal-mean training path: the MSTNet blocks in Figure~\ref{fig:mstnet-arch} are unchanged, and only the input aggregation differs, as in Figure~\ref{fig:data-flow-pixel}.

\section{Municipal-mean LOO diagnostics}
\label{app:muni-mean-eval}
Section~\ref{sec:exp-muni-mean} reports full-season municipal-mean LOO metrics.
Here we add the municipal-mean counterparts of the partial-season curves, scatterplots, Paraná error maps, and modality ablations from Section~\ref{sec:results}.
There is no Mambor\'{e} panel for this path because municipal-mean inputs do not produce intramunicipal scores.

\subsection{Partial-season evaluation}
Table~\ref{tab:incomplete-r2-muni} and Figure~\ref{fig:incomplete-r2-muni} report municipal-mean $R^2$ when only the first $k$ season months are kept.
As on the pixel path, early truncations are strongly negative; excluding 2022, mean $R^2$ reaches $-0.034$, $0.269$, and $0.275$ at four, five, and six months.

\begin{table}[!htbp]
\centering
\caption{Municipal-mean year-LOO $R^2$ under incomplete seasons.
Each column keeps only the first $k$ months of the Oct--Mar window ($k=6$ matches Table~\ref{tab:loo-muni-mean}).}
\label{tab:incomplete-r2-muni}
\begin{tabular}{lcccccc}
\toprule
LOO year & Oct & Nov & Dec & Jan & Feb & Mar \\
\midrule
2020 & $-19.078$ & $-1.517$ & $0.127$ & $0.072$ & $0.221$ & $0.233$ \\
2021 & $-19.950$ & $-18.097$ & $-1.441$ & $-0.605$ & $0.377$ & $0.410$ \\
2022 & $-0.975$ & $-0.858$ & $-0.524$ & $-0.360$ & $-0.308$ & $-0.300$ \\
2023 & $-3.056$ & $-4.617$ & $-0.643$ & $0.128$ & $0.220$ & $0.199$ \\
2024 & $-4.293$ & $-3.539$ & $0.253$ & $0.271$ & $0.256$ & $0.259$ \\
\midrule
Mean & $-9.470$ & $-5.725$ & $-0.446$ & $-0.099$ & $0.153$ & $0.160$ \\
Mean excl.\ 2022 & $-11.594$ & $-6.942$ & $-0.426$ & $-0.034$ & $0.269$ & $0.275$ \\
\bottomrule
\end{tabular}
\end{table}

\begin{figure}[!htbp]
\centering
\includegraphics[width=0.85\textwidth]{incomplete_series_r2_muni_mean.pdf}
\caption{Municipal-mean LOO $R^2$ versus season length.
Each curve is one left-out year; the black squares are the mean over non-drought folds.
Points below the axis floor are clipped and marked with downward triangles.}
\label{fig:incomplete-r2-muni}
\end{figure}

\subsection{Scatterplots and spatial maps}
Figure~\ref{fig:scatter-muni} shows predicted versus observed municipal productivity under the municipal-mean path for every LOO year.
Figure~\ref{fig:map-error-muni} shows municipal error maps for the best and worst LOO years; Figure~\ref{fig:map-error-muni-app} shows 2020, 2023, and 2024.

\begin{figure}[!htbp]
\centering
\includegraphics[width=\linewidth,height=0.42\textheight,keepaspectratio]{scatter_panel_muni_mean.pdf}
\caption{Predicted versus observed municipal soybean productivity (\si{t/ha}) for each year-LOO fold under the municipal-mean path.
Every panel evaluates only the left-out year. The dashed line is identity.}
\label{fig:scatter-muni}
\end{figure}

\begin{figure}[!htbp]
\centering
\includegraphics[width=\linewidth,keepaspectratio]{map_error_panel_muni_mean.png}
\caption{Municipal aggregate prediction error maps under the municipal-mean path for LOO years 2021 and 2022.}
\label{fig:map-error-muni}
\end{figure}

\begin{figure}[!htbp]
\centering
\includegraphics[width=\linewidth,keepaspectratio]{map_error_panel_muni_mean_appendix.png}
\caption{Municipal aggregate prediction error maps under the municipal-mean path for LOO years 2020, 2023, and 2024.}
\label{fig:map-error-muni-app}
\end{figure}

\subsection{Modality ablations}
\label{app:ablations-muni}
Table~\ref{tab:ablations-muni} reports the same leave-one-modality-out winners as Table~\ref{tab:ablations}, scored on the municipal-mean path used in training.
Excluding 2022, full MSTNet has mean $R^2=0.275$.
Omitting climate costs about $0.10$ in mean $R^2$, omitting Sentinel-2 about $0.07$, and omitting soil about $0.04$.

\begin{table}[!htbp]
\centering
\small
\caption{Year-LOO municipal-mean metrics for full MSTNet and leave-one-modality-out ablations.
Each ablation is the best of three architectures under municipal-mean year-LOO validation.
RMSE is in \si{t/ha}.
The first block is the full three-stream baseline; blocks below omit one modality.}
\label{tab:ablations-muni}
\begin{tabular}{lllccccccc}
\toprule
Variant & Inputs & Metric & 2020 & 2021 & 2022 & 2023 & 2024 & Mean & Mean excl.\ 2022 \\
\midrule
\multirow{2}{*}{Full MSTNet} & \multirow{2}{*}{S2+clim.+soil}
  & $R^2$ & 0.233 & 0.410 & $-0.300$ & 0.199 & 0.259 & 0.160 & 0.275 \\
  & & RMSE & 0.411 & 0.320 & 1.342 & 0.433 & 0.577 & 0.617 & 0.435 \\
\midrule
\multirow{2}{*}{No soil} & \multirow{2}{*}{S2+clim.}
  & $R^2$ & 0.207 & 0.303 & 0.076 & 0.133 & 0.295 & 0.203 & 0.234 \\
  & & RMSE & 0.418 & 0.348 & 1.132 & 0.450 & 0.563 & 0.582 & 0.445 \\
\midrule
\multirow{2}{*}{No climate} & \multirow{2}{*}{S2+soil}
  & $R^2$ & 0.292 & 0.233 & $-0.197$ & 0.145 & 0.035 & 0.101 & 0.176 \\
  & & RMSE & 0.394 & 0.365 & 1.288 & 0.447 & 0.658 & 0.630 & 0.466 \\
\midrule
\multirow{2}{*}{No Sentinel-2} & \multirow{2}{*}{clim.+soil}
  & $R^2$ & 0.131 & 0.262 & $-0.064$ & 0.167 & 0.258 & 0.151 & 0.204 \\
  & & RMSE & 0.438 & 0.358 & 1.214 & 0.441 & 0.578 & 0.606 & 0.454 \\
\bottomrule
\end{tabular}
\end{table}

\section{80/20 municipality holdout}
\label{app:muni80-20}
Section~\ref{sec:discussion-yoo-vs-muni} contrasts year-LOO with the municipality holdout used by~\citet{song2022soy}: randomly hold out 20\% of municipalities for testing, keep the other 80\% for training, and include every harvest year in both pools.
Table~\ref{tab:muni80-20} reports that holdout for the same Paran\'{a} MSTNet architecture as Table~\ref{tab:loo-muni-mean} (one Transformer layer, $d=128$, four heads, $d_{\mathrm{inner}}=64$, dropout $0.4$), trained once on the municipal-mean path with the 80\% training municipalities and scored on the held-out 20\%.
Per-year rows restrict the holdout to a single harvest year; Mean averages those five yearly scores; Pooled scores all holdout municipality--years together.
Pooled $R^2=0.736$ (RMSE $0.448$\,\si{t/ha}) is much higher than year-LOO on the same architecture, while the yearly $R^2$ values ($0.283$--$0.636$) show that much of the pooled score comes from shared seasonal structure across years rather than pure spatial interpolation within a fixed season.

\begin{table}[!htbp]
\centering
\caption{80/20 municipality holdout for Paran\'{a} MSTNet under the municipal-mean path.
Twenty percent of municipalities are held out for validation and test; all harvest years 2020--2024 appear in both pools.
RMSE and MAE are in \si{t/ha}.
Mean is the unweighted average of the five yearly scores; Pooled aggregates all holdout municipality--years.}
\label{tab:muni80-20}
\begin{tabular}{lcccc}
\toprule
Subset & $n$ & RMSE & MAE & $R^2$ \\
\midrule
2020 & 70 & 0.330 & 0.262 & 0.426 \\
2021 & 68 & 0.305 & 0.233 & 0.409 \\
2022 & 69 & 0.694 & 0.468 & 0.636 \\
2023 & 68 & 0.411 & 0.323 & 0.283 \\
2024 & 68 & 0.385 & 0.300 & 0.581 \\
\midrule
Mean & --- & 0.425 & 0.317 & 0.467 \\
Pooled & 343 & 0.448 & 0.317 & 0.736 \\
\bottomrule
\end{tabular}
\end{table}

\section{Additional pixel-path heatmaps}
\label{app:extra-heatmaps}
Section~\ref{sec:results} highlights the best and worst LOO years for the pixel-path Paraná and Mambor\'{e} maps.
Figures~\ref{fig:map-error-app} and~\ref{fig:mambore-app} report the remaining folds, 2020, 2023, and 2024.
Figure~\ref{fig:mambore-ndvi} stacks the 2022 Mambor\'{e} productivity map with mean NDVI and peak NDVI on the same soy pixels; we use 2022 because its peak NDVI field is less homogeneous than in the other LOO years, so intramunicipal contrast is easier to see.
That comparison is qualitative only.

\begin{figure}[!htbp]
\centering
\includegraphics[width=\linewidth,keepaspectratio]{map_error_panel_appendix.png}
\caption{Municipal aggregate prediction error maps under the pixel path for LOO years 2020, 2023, and 2024.}
\label{fig:map-error-app}
\end{figure}

\begin{figure}[!htbp]
\centering
\includegraphics[width=\linewidth,keepaspectratio]{mambore_panel_appendix.png}
\caption{Intramunicipal Mambor\'{e} productivity maps under the pixel path for LOO years 2020, 2023, and 2024.}
\label{fig:mambore-app}
\end{figure}

\begin{figure}[!htbp]
\centering
\includegraphics[width=0.92\linewidth,height=0.88\textheight,keepaspectratio]{mambore_ndvi_compare.png}
\caption{Mambor\'{e} under the 2022 LOO checkpoint: predicted productivity, mean NDVI, and peak NDVI on soy pixels, stacked so that spatial patterns can be compared directly.}
\label{fig:mambore-ndvi}
\end{figure}
